\section{Alliance Clusters}
\label{app:alliance}

The skills learned on Pelee and Rondeau transfer directly to the national HPC systems of the
Digital Research Alliance of Canada (the Alliance). To use them you need an Alliance account
(\url{https://ccdb.alliancecan.ca}) and multi-factor authentication. You do \emph{not} need the
university VPN to reach them.

\begin{table}[h]
\centering\small
\begin{tabular}{@{}lp{0.30\textwidth}cl@{}}
\toprule
\textbf{Cluster} & \textbf{Location} & \textbf{GPUs/node} & \textbf{Login node} \\
\midrule
Nibi     & University of Waterloo (SHARCNET)         & 8 & \texttt{nibi.alliancecan.ca} \\
Fir      & Simon Fraser University                   & 4 & \texttt{fir.alliancecan.ca} \\
Rorqual  & École de technologie supérieure, Montréal & 4 & \texttt{rorqual.alliancecan.ca} \\
Trillium & University of Toronto (SciNet)            & 4 & \texttt{trillium-gpu.alliancecan.ca} \\
\bottomrule
\end{tabular}
\caption{GPU nodes (NVIDIA H100, 80~GB) on four Alliance clusters. Check the Alliance
documentation for current details. Trillium's CPU nodes use a different login node
(\texttt{trillium.alliancecan.ca}), and the cluster is run differently from the others, so read
its documentation before your first job.}
\end{table}

Daily tasks remain similar:

\begin{lstlisting}
ssh cd ls tmux git
\end{lstlisting}

The major difference is that large jobs must be submitted through a scheduler (Slurm).